\thispagestyle{empty}
% \renewcommand{\chapterheadstartvskip}{\vspace*{-\topskip}}
% \renewcommand{\chapterheadendvskip}{%
%   \vspace*{1\baselineskip plus .1\baselineskip minus .167\baselineskip}}
\phantomsection
\addcontentsline{toc}{chapter}{Abbildungsverzeichnis}\label{cha:Abbildungsverzeichnis}
\listoffigures   
\section{Bild Quellen}
Einige Bilder, das Titelbild eingeschlossen, sind mit dezgo.com generiert worden. 
Dezgo ist eine free to use image generator KI. Eine Lupe, die ein Neuronales Netz umgeben von Computercode herausdeutet
und mit einem Warnhinweisschild versehen ist, wurde gewählt um KI Sicherheit darzustellen.
Ein großes Flugzeug im Hintergrund wurde gewählt um den Anwendungsfall in der der Luftfahrt zu symbolisieren 
und eine große Satellitenschüssel um die Anwendung für Radarsysteme zu verdeutlichen.
\paragraph{Part II Theoretische und technische Grundlagen}
\begin{itemize}
 \item chapter banner:
     title   = {Professor der komplexe Formeln an eine Tafel schreibt},
      author  = {{Unsplash contributor unknown}},   % optional: echten Namen einsetzen, falls angegeben
      url     = {https://unsplash.com/de/fotos/professor-der-komplexe-formeln-an-eine-tafel-schreibt-d4fGiVvfyPk},
      urldate = {2025-10-03},
      note    = {Unsplash}
\end{itemize}

\newpage

%\renewcommand{\chapterheadstartvskip}{\vspace*{-\topskip}}
%\renewcommand{\chapterheadendvskip}{%
%  \vspace*{1\baselineskip plus .1\baselineskip minus .167\baselineskip}}